\documentclass[11pt]{amsart}

\usepackage{amsmath}
\usepackage{amssymb}
\usepackage{graphicx}

\title{Mathematical Theory Behind The Semantic Turing Field}
\author{Karim El-Sharkawy}

\textwidth 5.50in
\oddsidemargin 0.375in
\evensidemargin 0.375in

\begin{document}

\maketitle

\section{Introduction}

The Semantic Turing Field (STF) is a computational model in which words are
treated as particles whose interactions are governed by semantic relationships.
The motivation is to combine a pretrained semantic representation of language
with a dynamical system: a word embedding assigns each word a vector in a
high-dimensional semantic space, and pairwise similarity in that space is used
to decide whether two particles attract or repel one another. The resulting
forces then evolve a separate two-dimensional particle field over time.

It is important to keep two distinct spaces in mind throughout. The
\emph{semantic space} is where the embeddings live and where relationships
between words are encoded. The \emph{simulation space} is where particles
actually move. The purpose of the simulation is not to reproduce the geometry
of the embedding space exactly, but rather to provide a dynamic visualization
of the semantic relationships that the embedding encodes.

\section{Semantic Representation}

Let the vocabulary consist of $N$ words. Each word $i$ is represented by an
embedding vector $\mathbf{v}_i \in \mathbb{R}^d$, where $d$ is the embedding
dimension. In the current implementation these vectors are obtained from GloVe.

Semantic similarity between two words is measured by cosine similarity,
\begin{equation}
\label{eq:cosine}
s_{ij}
=
\frac{\mathbf{v}_i \cdot \mathbf{v}_j}{\|\mathbf{v}_i\|\,\|\mathbf{v}_j\|},
\end{equation}
and the collection of all pairwise similarities forms the matrix
$S = (s_{ij}) \in \mathbb{R}^{N \times N}$. This matrix is the sole source of
semantic information used by the force model; everything that follows is a
consequence of it.

\section{Simulation Space}

Each word is associated with a dynamic particle position
$\mathbf{x}_i(t) \in \mathbb{R}^2$ and a velocity
$\mathbf{u}_i(t) \in \mathbb{R}^2$. The particle position is \emph{not} the
original GloVe embedding: the embeddings are first reduced to two dimensions to
provide an initial configuration, after which the simulation evolves the
particle positions dynamically. The state of the field at time $t$ is therefore
the collection
\[
\mathcal{S}(t) = \{\mathbf{x}_i(t), \mathbf{u}_i(t)\}_{i=1}^{N}.
\]
Throughout a normal simulation the semantic information is held fixed, while
the particle state evolves.

\section{The Pairwise Force Model}

The fundamental interaction between two particles is determined by their
semantic similarity. Writing $\mathbf{r}_{ij} = \mathbf{x}_j - \mathbf{x}_i$
for the displacement from particle $i$ toward particle $j$, the STF pairwise
force is
\begin{equation}
\label{eq:pairwise-force}
\boxed{
\mathbf{F}_{ij}
=
\alpha \left( s_{ij} - \beta \right) \mathbf{r}_{ij},
}
\end{equation}
where $s_{ij}$ is the semantic similarity of \eqref{eq:cosine}, $\alpha > 0$
controls the overall interaction strength, $\beta$ is a similarity threshold
separating attraction from repulsion, and $\mathbf{r}_{ij}$ gives the spatial
direction from $i$ toward $j$. The total pairwise force acting on particle $i$
is the sum over all other particles,
\begin{equation}
\label{eq:total-semantic}
\mathbf{F}_i^{\mathrm{semantic}}
=
\sum_{j \neq i} \mathbf{F}_{ij}.
\end{equation}

\subsection{Attraction and Repulsion}

The qualitative behaviour of the interaction is controlled entirely by the sign
of $s_{ij} - \beta$. If $s_{ij} > \beta$ then the coefficient is positive and
the force points toward particle $j$, so the interaction is attractive. If
$s_{ij} < \beta$ then the coefficient is negative and the force reverses
direction, so the interaction is repulsive. When $s_{ij} = \beta$ the pairwise
interaction vanishes. This is summarized in Table~\ref{tab:interaction}.

\begin{table}[h]
\centering
\begin{tabular}{c|c}
\textbf{Condition} & \textbf{Interaction} \\
\hline
$s_{ij} > \beta$ & Attraction \\
$s_{ij} = \beta$ & No pairwise interaction \\
$s_{ij} < \beta$ & Repulsion \\
\end{tabular}
\caption{The sign of $s_{ij} - \beta$ determines the qualitative interaction.}
\label{tab:interaction}
\end{table}

The threshold $\beta$ is therefore not merely a scaling parameter: it
determines which semantic relationships are treated as cohesive and which as
separating.

\subsection{Why Both Attraction and Repulsion Are Needed}

It is worth pausing to see why the signed interaction is essential. If every
pair of particles attracted one another, the field would collapse into a small
region. If every pair repelled, the particles would drift apart without forming
any recognizable structure. The signed force provides both local cohesion and
global separation at once. The resulting organization is an emergent property
of many simultaneous pairwise interactions rather than the output of any
explicit clustering procedure.

\section{Relationship to Spring Systems}

The pairwise force in \eqref{eq:pairwise-force} is structurally similar to a
spring or graph-based force, because the interaction acts along the line
connecting the two particles. A conventional Hookean spring has the form
$\mathbf{F} = -k(\mathbf{x}_i - \mathbf{x}_j)$, where the spatial displacement
alone determines the magnitude of the force. In STF, by contrast, the spatial
displacement determines only the direction, while semantic similarity
determines the strength and sign:
\[
\text{spatial displacement} \rightarrow \text{direction},
\qquad
\text{semantic similarity} \rightarrow \text{interaction strength and sign}.
\]
The model can therefore be viewed as a semantic analogue of a force-directed
system. The important distinction is that STF does not treat spatial distance
itself as the primary source of semantic information; the semantic relationship
originates in the embedding space.

\section{Energy Interpretation}

The pairwise interaction admits an energy interpretation. For a fixed
similarity value $s_{ij}$, define the pairwise potential
\begin{equation}
\label{eq:pairwise-energy}
E_{ij}
=
\frac{\alpha}{2} \left( s_{ij} - \beta \right)
\|\mathbf{x}_i - \mathbf{x}_j\|^2.
\end{equation}
Taking the negative gradient with respect to $\mathbf{x}_i$ gives
\[
-\nabla_{\mathbf{x}_i} E_{ij}
=
\alpha \left( s_{ij} - \beta \right) (\mathbf{x}_j - \mathbf{x}_i),
\]
which is exactly the STF pairwise force \eqref{eq:pairwise-force}, so that
$\mathbf{F}_{ij} = -\nabla_{\mathbf{x}_i} E_{ij}$.

The interpretation of this potential depends on the sign of $s_{ij} - \beta$.
When $s_{ij} > \beta$ the coefficient of the quadratic is positive, so the
energy is minimized when the two particles are close. When $s_{ij} < \beta$ the
coefficient is negative, so increasing separation lowers the pairwise energy;
in an unconstrained system this would encourage particles to move arbitrarily
far apart. This is one reason the actual simulation requires additional
mechanisms such as boundary forces and damping. The energy interpretation
should therefore be understood as a description of the pairwise semantic
interaction, not as a claim that the full implemented system is pure gradient
descent on a single global energy function.

\section{Discrete-Time Dynamics}

The simulation evolves according to a simple discrete-time integration scheme.
The total force on particle $i$ is
\begin{equation}
\label{eq:total-force}
\mathbf{F}_i(t)
=
\mathbf{F}_i^{\mathrm{semantic}}
+ \mathbf{F}_i^{\mathrm{gravity}}
+ \mathbf{F}_i^{\mathrm{boundary}},
\end{equation}
and the velocity and position are updated by
\begin{align}
\mathbf{u}_i(t + \Delta t)
&=
\mathbf{u}_i(t) + \mathbf{F}_i(t)\,\Delta t, \label{eq:velocity-update} \\
\mathbf{x}_i(t + \Delta t)
&=
\mathbf{x}_i(t) + \mathbf{u}_i(t + \Delta t)\,\Delta t. \label{eq:position-update}
\end{align}
Damping is applied after the integration step, $\mathbf{u}_i \leftarrow
\gamma \mathbf{u}_i$, where $\gamma$ is the damping parameter. The
implementation also limits the maximum particle speed in order to maintain
numerical stability and visual control.

\section{Boundary Forces}

The semantic interaction alone does not constrain particles to remain within
the visible simulation region, so a boundary force
$\mathbf{F}_i^{\mathrm{boundary}}$ is added to the total force
\eqref{eq:total-force}. Boundary forces should be interpreted as simulation
constraints rather than as semantic relationships.

\section{K-Means Clustering}

K-Means is used as a complementary structural analysis of the embedding space.
Given the word embeddings $\{\mathbf{v}_1, \ldots, \mathbf{v}_N\}$, K-Means
partitions them into $K$ clusters by minimizing
\[
\sum_{k=1}^{K} \sum_{\mathbf{v}_i \in C_k}
\|\mathbf{v}_i - \boldsymbol{\mu}_k\|^2,
\]
where $\boldsymbol{\mu}_k$ is the centroid of cluster $C_k$. The resulting
cluster assignment for word $i$ is denoted $c_i \in \{1, \ldots, K\}$.

K-Means does not determine the physical force between particles. Cosine
similarity drives the semantic interaction, while K-Means supplies cluster
assignments for visualization. This distinction matters because a pair of words
can belong to the same K-Means cluster without necessarily having a strong
direct interaction under the STF force model. K-Means therefore provides a
coarse semantic grouping, whereas the force field represents continuous
pairwise relationships.

\section{Gravity Wave}

The gravity wave introduces an external semantic perturbation into the field.
Suppose a user enters a sentence $q$. After preprocessing, the sentence is
represented by the mean embedding
\[
\mathbf{q}
=
\frac{1}{|V_q|} \sum_{w \in V_q} \mathbf{v}_w,
\]
where $V_q$ is the set of valid vocabulary words appearing in the sentence. For
every vocabulary word $i$ we compute the similarity
$s_{iq} = \operatorname{cosine\_similarity}(\mathbf{v}_i, \mathbf{q})$, and the
gravity-wave force is activated for sufficiently relevant words. A simplified
representation of this force is
\begin{equation}
\label{eq:gravity}
\mathbf{F}_i^{\mathrm{gravity}}
=
-\lambda \max(s_{iq} - \tau, 0)\, \mathbf{x}_i,
\end{equation}
where $\lambda$ controls the strength of the external influence, $\tau$ is the
relevance threshold, and $s_{iq}$ measures similarity between word $i$ and the
input sentence. Because the force points toward the origin, words semantically
related to the input sentence are drawn toward the center of the field. The
gravity wave thus acts as a semantic query mechanism:
\[
\text{sentence} \rightarrow \text{sentence embedding}
\rightarrow \text{word--sentence similarity}
\rightarrow \text{external force}
\rightarrow \text{field response}.
\]
Unlike the pairwise semantic interaction, the gravity wave is externally
triggered and temporary.

\section{Emergent Structure}

No explicit rule instructs the simulation to form semantic communities;
structure emerges from the interaction of many particles. The overall process
can be summarized as
\[
\begin{gathered}
\text{embeddings}\\
\downarrow\\
\text{similarity matrix}\\
\downarrow\\
\text{signed interactions}\\
\downarrow\\
\text{particle dynamics}\\
\downarrow\\
\text{emergent spatial organization}
\end{gathered}
\]
Semantically related words tend to experience attractive interactions when
their similarity exceeds the threshold, while less-related words experience
repulsive interactions. The resulting spatial organization is therefore a
dynamical representation of the semantic relationships encoded by the embedding
model.

\section{What STF Represents}

STF should be interpreted as a \emph{semantic dynamical visualization} rather
than as a replacement for the original embedding model. The embedding model
supplies the semantic information; the STF force model turns that information
into interactions; the simulation integrates those interactions over time;
K-Means provides a coarse structural grouping for visual interpretation; and
the renderer exposes the resulting system to the user. The complete conceptual
pipeline is
\[
\boxed{
\text{language} \rightarrow \text{embedding} \rightarrow \text{similarity}
\rightarrow \text{force field} \rightarrow \text{dynamics}
\rightarrow \text{visualization},
}
\]
with K-Means providing an additional clustering layer and the gravity wave
providing an interactive semantic perturbation.

\section{Limitations of the Current Model}

The current STF formulation is intentionally simple, and several limitations
should be recognized. First, the semantic representation depends entirely on
the pretrained embedding model, so any biases or limitations in the embeddings
are inherited by the field. Second, the two-dimensional simulation space is not
itself a faithful representation of the original high-dimensional semantic
space; it is a dynamic visualization space. Third, the threshold $\beta$ is a
manually selected parameter, and different values can produce substantially
different field structures. Fourth, the current force model does not explicitly
include mass, realistic physical units, or a physically derived potential
governing every component of the system. Finally, K-Means provides a useful
visualization aid but should not be interpreted as ground-truth semantic
categories. These limitations are intentional in the current version: STF is
primarily a computational framework for exploring how semantic relationships
can be represented dynamically.

\section{Future Extensions}

The current formulation provides a foundation for more sophisticated
experiments. Natural extensions include alternative embedding models and
clustering algorithms; normalized or distance-dependent force magnitudes;
additional force components; spatial partitioning for larger vocabularies;
quantitative measures of field stability; systematic experiments over
$\alpha$, $\beta$, $\Delta t$, and the damping parameter; more principled
energy formulations; interactive semantic queries; and reproducible simulation
and animation generation. The architecture is intentionally modular so that
these extensions can be introduced without replacing the entire system.

\end{document}